% Part: normal-modal-logic
% Chapter: tableaux
% Section: rules-for-K

\documentclass[../../../include/open-logic-section]{subfiles}

\begin{document}

\olfileid{nml}{tab}{rul}

\olsection{د \Ax{K} لپاره قاعدې}

د عامو قضيېيزو رابطو قاعدې د عامو نښه لرونکو قضيېيزو
!!{tableau}s له قاعدو سره يو شان دي؛ يوازې مخکښونه ورسره
زياتېږي۔ په هر حالت کښې، پر نښه لرونکي !!{formula}
$\sFmla{S}{!A}[\sigma]$ پلې شوې قاعده نوي !!{formula}s
جوړوي چې د هغو مخکښ هم~$\sigma$ وي۔ دا بايد له مفهومي پلوه
څرګنده وي: که $!A \land !B$ په~$\sigma$ نومې نړۍ کښې رښتينے
وي، نو $!A$ او $!B$ هم په~$\sigma$ کښې رښتيني دي (نۀ په
بله نړۍ کښې)۔ قضيېيزې قاعدې په \olref{tab:prop-rules} کښې
راټولې شوې دي۔

\begin{table}
  \[\def\arraystretch{3}\begin{array}{|c|c|}
    \hline
    \AxiomC{\sFmla{\True}{\lnot !A}[\sigma]}
    \RightLabel{\TRule{\True}{\lnot}}
    \UnaryInfC{\sFmla{\False}{!A}[\sigma]}
    \DisplayProof
    &
    \AxiomC{\sFmla{\False}{\lnot !A}[\sigma]}
    \RightLabel{\TRule{\False}{\lnot}}
    \UnaryInfC{\sFmla{\True}{!A}[\sigma]}
    \DisplayProof
    \\[1ex]
    \hline
    \AxiomC{\sFmla{\True}{!A \land !B}[\sigma]}
    \RightLabel{\TRule{\True}{\land}}
    \UnaryInfC{\sFmla{\True}{!A}[\sigma]}
    \noLine
    \UnaryInfC{\sFmla{\True}{!B}[\sigma]}
    \DisplayProof
    &
    \AxiomC{\sFmla{\False}{!A \land !B}[\sigma]}
    \RightLabel{\TRule{\False}{\land}}
    \UnaryInfC{$\sFmla{\False}{!A}[\sigma] \quad \mid \quad
      \sFmla{\False}{!B}[\sigma]$}
    \DisplayProof
    \\[2ex]
    \hline
    \AxiomC{\sFmla{\True}{!A \lor !B}[\sigma]}
    \RightLabel{\TRule{\True}{\lor}}
    \UnaryInfC{$\sFmla{\True}{!A}[\sigma] \quad \mid \quad
      \sFmla{\True}{!B}[\sigma]$}
    \DisplayProof
    &
    \AxiomC{\sFmla{\False}{!A \lor !B}[\sigma]}
    \RightLabel{\TRule{\False}{\lor}}
    \UnaryInfC{\sFmla{\False}{!A}[\sigma]}
    \noLine
    \UnaryInfC{\sFmla{\False}{!B}[\sigma]}
    \DisplayProof
    \\[2ex]
    \hline
    \AxiomC{\sFmla{\True}{!A \lif !B}[\sigma]}
    \RightLabel{\TRule{\True}{\lif}}
    \UnaryInfC{$\sFmla{\False}{!A}[\sigma] \quad \mid
      \quad \sFmla{\True}{!B}[\sigma]$}
    \DisplayProof
    &
    \AxiomC{\sFmla{\False}{!A \lif !B}[\sigma]}
    \RightLabel{\TRule{\False}{\lif}}
    \UnaryInfC{\sFmla{\True}{!A}[\sigma]}
    \noLine
    \UnaryInfC{\sFmla{\False}{!B}[\sigma]}
    \DisplayProof
    \\[2ex]
    \hline
  \end{array}\]
  \caption{د قضيېيزو رابطو لپاره د مخکښ‌لرونکي !!{tableau} قاعدې}
  \ollabel{tab:prop-rules}
\end{table}

د تړلو شرط د عامو !!{tableau}s د شرط په شان دے، خو دلته بايد
يوازې !!{formula}s نه، بلکې مخکښونه هم سره يو شان وي۔ نو که
په يوه څانګه کښې دواړه
\[
\sFmla{\True}{!A}[\sigma] \quad\text{او}\quad \sFmla{\False}{!A}[\sigma]
\]
د کوم مخکښ~$\sigma$ او !!{formula}~$!A$ لپاره وي، هغه څانګه
تړلې ده۔

د فرضونو د ايښودلو قاعدې هم د عامو !!{tableau}s په شان دي، خو
د فرضونو لپاره تل مخکښ~$1$ کاروو۔ (د مخکښ ټاکنه مهمه نۀ ده،
تر هغو چې د ټولو فرضونو لپاره يو شان وي۔) نو، د بېلګې په توګه،
\[
!B_1, \dots, !B_n \Proves !A
\]
هله او يوازې هله وايو چې د لاندې فرضونو لپاره تړلے تابلو وي:
\[
\sFmla{\True}{!B_1}[1], \dots, \sFmla{\True}{!B_n}[1],
\sFmla{\False}{!A}[1].
\]

د مودالي \iftag{prvBox}{\iftag{prvDiamond}{نښو~$\Box$
    او}{نښې~$\Box$}}{نښې}\iftag{prvDiamond}{~$\Diamond$}{}
لپاره، پر هغه !!a{formula} چې مخکښ يې~$\sigma$ دے د قاعدې
د پلي کېدو پايلې مخکښ $\sigma.n$ وي۔ خو کوم $n$ روا دے، دا
په دې پورې تړلې ده چې نښه~$\True$ ده که~$\False$۔

\iftag{prvBox}{د $\TRule{\True}{\Box}$ قاعده هغه څانګه چې
  $\sFmla{\True}{\Box !A}[\sigma]$ لري، په
  $\sFmla{\True}{!A}[\sigma.n]$ غځوي۔\iftag{prvDiamond}{
    همداسې، }{}}{}\iftag{prvDiamond}{د
  $\TRule{\False}{\Diamond}$ قاعده هغه څانګه چې
  $\sFmla{\False}{\Diamond !A}[\sigma]$ لري، په
  $\sFmla{\False}{!A}[\sigma.n]$ غځوي۔}{}
\iftag{notprvBox,notprvDiamond}{دا قاعده}{دا قاعدې} يوازې د
هغې~$\sigma.n$ لپاره پلې کېدے شي چې په اړونده څانګه کښې
\emph{لا مخکې} راغلې وي۔ داسې مخکښ په څانګه کښې «کارېدلے» بولو۔

\iftag{prvBox}{د $\TRule{\False}{\Box}$ قاعده هغه څانګه چې
  $\sFmla{\False}{\Box !A}[\sigma]$ لري، په
  $\sFmla{\False}{!A}[\sigma.n]$ غځوي۔\iftag{prvDiamond}{
    همداسې، }{}}{}\iftag{prvDiamond}{د
  $\TRule{\True}{\Diamond}$ قاعده هغه څانګه چې
  $\sFmla{\True}{\Diamond !A}[\sigma]$ لري، په
  $\sFmla{\True}{!A}[\sigma.n]$ غځوي۔}{}
خو \iftag{notprvBox,notprvDiamond}{دا قاعده}{دا قاعدې} يوازې
د هغې~$\sigma.n$ لپاره پلې کېدے شي چې په اړونده څانګه کښې
\emph{لا مخکې نۀ} وي راغلې۔ داسې مخکښ په څانګه کښې «نوے» بولو۔

قاعدې په \olref{tab:rules-K} کښې ورکړل شوې دي۔

\begin{table}
  \begin{center}
    \def\arraystretch{3}\def\fCenter{}
    \begin{tabular}{|c|c|}
    \hline
    \iftag{prvBox}{
      \AxiomC{\sFmla{\True}{\Box !A}[\sigma]}
      \RightLabel{\TRule{\True}{\Box}}
      \UnaryInfC{\sFmla{\True}{!A}[\sigma.n]}
      \DisplayProof
      &
      \AxiomC{\sFmla{\False}{\Box !A}[\sigma]}
      \RightLabel{\TRule{\False}{\Box}}
      \UnaryInfC{\sFmla{\False}{!A}[\sigma.n]}
      \DisplayProof\\
      $\sigma.n$ کارېدلے دے & $\sigma.n$ نوے دے
      \\[1ex]
      \hline}{}
    \iftag{prvDiamond}{
      \AxiomC{\sFmla{\True}{\Diamond !A}[\sigma]}
      \RightLabel{\TRule{\True}{\Diamond}}
      \UnaryInfC{\sFmla{\True}{!A}[\sigma.n]}
      \DisplayProof
      &
      \AxiomC{\sFmla{\False}{\Diamond !A}[\sigma]}
      \RightLabel{\TRule{\False}{\Diamond}}
      \UnaryInfC{\sFmla{\False}{!A}[\sigma.n]}
      \DisplayProof\\
      $\sigma.n$ نوے دے & $\sigma.n$ کارېدلے دے
      \\[1ex]
      \hline}{}
    \end{tabular}
  \end{center}
  \caption{د \Ax{K} مودالي قاعدې۔}
  \ollabel{tab:rules-K}
\end{table}

د \iftag{prvBox}{\TRule{\True}{\Box}}{\TRule{\False}{\Diamond}}
لپاره د مخکښ د مخکې کارېدلو قيد اړين دے؛ ګنې لاندې تابلو به
هم تړلے !!{tableau} ګڼل کېده:
\iftag{notprvBox,notprvDiamond}{%
  \iftag{prvBox}{%
  \begin{oltableau}
    [\pFmla{\True}{\Box \formula{A}}{1}, just = \TAss
      [\pFmla{\False}{\lnot\Box\lnot \formula{A}}{1}, just = \TAss
        [\pFmla{\True}{\formula{A}}{1.1}, just = {\TRule{\True}{\Box}[1]}
          [\pFmla{\True}{\Box\lnot\formula{A}}{1},
            just ={\TRule{\False}{\lnot}[2]}
            [\pFmla{\True}{\lnot\formula{A}}{1.1},
              just = {\TRule{\True}{\Box}[4]}
              [\pFmla{\False}{\formula{A}}{1.1},
                  just ={\TRule{\True}{\lnot}[5]}, close]
            ]
          ]
        ]
      ]
    ]
  \end{oltableau}
  }{
  \begin{oltableau}
    [\pFmla{\True}{\lnot\Diamond\lnot \formula{A}}{1}, just = \TAss
      [\pFmla{\False}{\Diamond\formula{A}}{1}, just = \TAss
        [\pFmla{\False}{\formula{A}}{1.1}, just = {\TRule{\False}{\Diamond}[2]}
          [\pFmla{\False}{\Diamond\lnot\formula{A}}{1},
            just ={\TRule{\True}{\lnot}[1]}
            [\pFmla{\False}{\lnot\formula{A}}{1.1},
              just = {\TRule{\False}{\Diamond}[4]}
              [\pFmla{\True}{\formula{A}}{1.1},
                  just ={\TRule{\True}{\lnot}[5]}, close]
            ]
          ]
        ]
      ]
    ]
  \end{oltableau}
}}{
  \begin{oltableau}
    [\pFmla{\True}{\Box \formula{A}}{1}, just = \TAss
      [\pFmla{\False}{\Diamond \formula{A}}{1}, just = \TAss
        [\pFmla{\True}{\formula{A}}{1.1}, just = {\TRule{\True}{\Box}[1]}
          [\pFmla{\False}{\formula{A}}{1.1},
            just = {\TRule{\False}{\Diamond}[2]}, close]
        ]
      ]
    ]
\end{oltableau}}

خو $\Box \formula{A} \Entails/ \Diamond \formula{A}$؛ نو زموږ
ثبوتي نظام به ناسم وي۔ همداسې، $\Diamond \formula{A} \Entails/
\Box \formula{A}$؛ خو که د
\iftag{prvBox}{\TRule{\False}{\Box}}{\TRule{\True}{\Diamond}}
لپاره د نوي مخکښ قيد نۀ وي، لاندې تابلو به تړلے وي:
\iftag{notprvBox,notprvDiamond}{%
  \iftag{prvBox}{%
  \begin{oltableau}
    [\pFmla{\True}{\lnot\Box\lnot \formula{A}}{1}, just = \TAss
      [\pFmla{\False}{\Box \formula{A}}{1}, just = \TAss
        [\pFmla{\False}{\formula{A}}{1.1}, just = {\TRule{\True}{\Box}[2]}
          [\pFmla{\False}{\Box\lnot\formula{A}}{1},
            just ={\TRule{\True}{\lnot}[1]}
            [\pFmla{\False}{\lnot\formula{A}}{1.1},
              just = {\TRule{\False}{\Box}[4]}
              [\pFmla{\True}{\formula{A}}{1.1},
                  just ={\TRule{\False}{\lnot}[5]}, close]
            ]
          ]
        ]
      ]
    ]
  \end{oltableau}
  }{
  \begin{oltableau}
    [\pFmla{\True}{\Diamond \formula{A}}{1}, just = \TAss
      [\pFmla{\False}{\lnot\Diamond\lnot\formula{A}}{1}, just = \TAss
        [\pFmla{\True}{\formula{A}}{1.1}, just = {\TRule{\True}{\Diamond}[1]}
          [\pFmla{\True}{\Diamond\lnot\formula{A}}{1},
            just ={\TRule{\False}{\lnot}[2]}
            [\pFmla{\True}{\lnot\formula{A}}{1.1},
              just = {\TRule{\True}{\Diamond}[4]}
              [\pFmla{\False}{\formula{A}}{1.1},
                  just ={\TRule{\True}{\lnot}[5]}, close]
            ]
          ]
        ]
      ]
    ]
  \end{oltableau}
}}{
  \begin{oltableau}
    [\pFmla{\True}{\Diamond \formula{A}}{1}, just = \TAss,name=one
      [\pFmla{\False}{\Box \formula{A}}{1}, just = \TAss, name=two
        [\pFmla{\True}{\formula{A}}{1.1}, just = {\TRule{\True}{\Diamond}[1]}
          [\pFmla{\False}{\formula{A}}{1.1},
            just = {\TRule{\False}{\Box}[2]}, close]
        ]
      ]
    ]
  \end{oltableau}}

\end{document}
